\documentclass{situationist-preprint}

\title{Photography, Transparency and Art}
\msid{AES-14-a}
\msfamily{AES-14}
\msdiscipline{Aesthetics}
\mstype{Principal manuscript}
\mskeywords{photography; transparency; Scruton; Walton; framing; selection; representation; artistic agency}

\begin{document}
\maketitle

\begin{abstract}
Is photography art? Scruton argued that it is not a representational art, because a photograph stands in a causal relation to its subject that leaves no room for the photographer's thought to fix what it shows, so that interest in a photograph is interest in its subject. Walton argued that photographs are transparent: in looking at a photograph we literally see, indirectly, what it is a photograph of. We accept transparency and reject Scruton's conclusion. Transparency is a counterfactual dependence of what a photograph shows on the scene, given the camera's position, frame, moment and settings; it says nothing about how those are chosen. We argue that the choices of viewpoint, frame and moment, and the selection of one exposure among many, determine which seeing a photograph offers, that they fix aesthetic properties of the photograph that the scene does not have, and that they are appreciable as the photographer's achievement. Photography is therefore art despite its causal transparency, and it is art through framing and selection, not through the accuracy that all photographs of a scene share.
\end{abstract}

\section{The problem}

From its beginning photography was described as the work of nature rather than of the artist. Fox Talbot called the photographic process the pencil of nature, and presented its images as impressed by light rather than drawn \citep{Talbot1844}. \citet{Bazin1960} held that the photographic image owes its power to an ontological fact: it is formed automatically, without the creative intervention of man, and so shares in the being of its model. The same fact that gave photography its documentary authority was taken to deny it the status of art. If an image is made by a mechanism, what is there in it for the maker to have done?

\citet{Scruton1981} gave the argument its sharpest form. A painting represents its subject through the painter's intention: what it shows is what the painter thought and chose to show, and its subject need not exist. An ideal photograph, by contrast, stands in a causal relation to its subject, which must exist and must have looked as the photograph shows it. Its content is fixed by the causal process, not by any thought of the photographer. It follows, Scruton argued, that the photograph is not a representation in the sense in which a painting is, and that any interest we take in it as a photograph is interest in what it shows: we look at it as a surrogate for its subject. Whatever artistic interest a photograph has derives from the respects in which the photographer has departed from the ideal, by staging, retouching or manipulation, and to that extent the result is a kind of painting.

\citet{Walton1984} defended a thesis about the nature of photographs that seems to strengthen Scruton's premise. Photographs are transparent: we see the world through them, as we see through mirrors, telescopes and spectacles. When I look at a photograph of my late grandfather, I see my grandfather, indirectly but literally. The ground of transparency is that the visual experience a photograph produces depends counterfactually on the scene in a way that does not pass through the beliefs of the photographer: had the scene been different, the photograph would have been different, whatever the photographer believed. A painting's dependence on its subject passes through the painter's beliefs, and so we see the subject in it but do not see the subject through it.

Responses have taken several forms. \citet{SnyderAllen1975} denied that photographs are mechanical records at all, arguing that every photograph embodies decisions about lens, exposure, development and printing. \citet{King1992} argued that Scruton's ideal photograph is not the kind of photograph that interests us. \citet{Currie1991} and \citet{CohenMeskin2004} examined the epistemic difference between photographs and handmade pictures, the latter arguing that it consists in photographs' independence of the maker's beliefs together with their supplying spatial information about the scene only relative to the viewer's background knowledge. \citet{Lopes2016} and \citet{Costello2018} have more recently distinguished the several arts that use photography and challenged the assumption that the photographic process is a single causal channel. \citet{Friday2002} gives a survey of the debate.

We take a different route from those who deny transparency. We grant Walton's thesis and argue that Scruton's conclusion does not follow from it. Transparency concerns the dependence of a photograph on its scene \emph{given} the camera's position, frame, moment and settings. It is silent about how these are fixed, and they are fixed by the photographer. Photography is art through them.

\section{Framework}

\begin{definition}[Transparency]\label{d:transp}
A picture is \emph{transparent} if the visual experience it produces depends counterfactually on the visible features of the scene before the device that made it, in a way that does not depend on the beliefs of any person about the scene \citep{Walton1984}.
\end{definition}

\begin{definition}[Photographic choices]\label{d:choices}
The \emph{framing} of a photograph consists of the choice of viewpoint, the boundaries of the frame, the moment of exposure, and the settings of the device that determine focus, depth of field, exposure and the rendering of tone and color. \emph{Selection} is the choice of one exposure, among those made, to present as the photograph, together with its cropping and its placement in a sequence or series.
\end{definition}

\begin{principle}[Art through agency]\label{p:art}
A practice of making objects in a medium is an art if, and only if, the objects it produces can have aesthetic properties and artistic value that are fixed by the maker's choices, and that are appreciable as the maker's achievement.
\end{principle}

\begin{principle}[Achievement]\label{p:ach}
A work's artistic value includes, besides its aesthetic value, the value of the achievement it embodies, assessed relative to the maker's situation (AES-04).
\end{principle}

Principle~\ref{p:art} states a sufficient condition that critics of photography and its defenders both accept: Scruton's argument is designed to show that photography fails it, not that it is the wrong test. It is compatible with the institutional account of AES-01, on which art is what is presented as art within an artworld subject to critical uptake: photography has long been presented and received as art, and the question is whether that uptake is warranted or, as Scruton held, rests on a confusion. Principle~\ref{p:ach} is the account of artistic value used in AES-04 and AES-16.

\section{The argument}

\begin{lemma}[Transparency is conditional]\label{l:cond}
Transparency is compatible with the dependence of a photograph's content on its maker's intentional choices of framing and selection.
\end{lemma}

\begin{proof}
The counterfactual of Definition~\ref{d:transp} holds fixed the device, its position and settings, and the moment of exposure: it says that, these being as they are, the photograph would have differed had the scene differed, whatever anyone believed about the scene. It does not say that the photograph would have been the same had these been different, and it would not have been. Change the viewpoint, the frame or the moment and the same scene yields a different photograph. Nor does transparency concern selection, which operates on the photographs that the process has produced. The photographer's beliefs and intentions are not in the causal chain from the scene to the image, given the framing; they are in the causal chain that fixes the framing and the selection. Both dependences hold together.
\end{proof}

\begin{lemma}[Framing fixes aesthetic properties]\label{l:aes}
The aesthetic properties of a photograph depend on its framing and selection, and are not properties of the scene.
\end{lemma}

\begin{proof}
Composition, balance, the tension between figure and edge, the rhythm of repeated forms within a frame, and the poise of a gesture caught at one instant rather than another are aesthetic properties of a photograph. Each is a property of a bounded, two-dimensional, instantaneous view, and the scene has no bounds, no picture plane and no privileged instant. The same street at the same hour yields one photograph whose figures are balanced against an empty doorway, and another, a step to the left and a second later, in which they are crowded and inert. The difference between them is a difference of aesthetic properties, and it is fixed by framing and selection, since the scene is the same.
\end{proof}

\begin{lemma}[Attributability]\label{l:attr}
The aesthetic properties fixed by framing and selection are appreciable as the photographer's achievement.
\end{lemma}

\begin{proof}
By Lemma~\ref{l:cond}, framing and selection are fixed by the photographer's choices, and by Lemma~\ref{l:aes} they fix aesthetic properties. An achievement is the bringing about of something of value through the exercise of competence, in a situation in which it could have failed. Securing a composition in a moving scene requires anticipation, perceptual discrimination and command of the device, as Cartier-Bresson's conception of the decisive moment, the instant at which the significance of an event and the organization of its forms coincide, makes explicit \citep{CartierBresson1952}; and selecting one exposure from many requires discrimination of which realizes those properties best. These competences can be exercised well or badly, and the result is assessable relative to what the photographer's situation allowed. By Principle~\ref{p:ach}, the properties so achieved contribute to the photograph's artistic value as the photographer's achievement.
\end{proof}

\begin{theorem}\label{t:main}
Photography is an art, despite its causal transparency, and it is an art through choices of framing and selection.
\end{theorem}

\begin{proof}
By Lemmas~\ref{l:aes} and~\ref{l:attr}, photographs can have aesthetic properties and artistic value fixed by the photographer's choices and appreciable as her achievement. By Principle~\ref{p:art}, photography is therefore an art. By Lemma~\ref{l:cond}, this is consistent with transparency, which we have not denied. The choices through which it is an art are framing and selection, since these are what Lemma~\ref{l:aes} shows to fix the relevant properties.
\end{proof}

\begin{corollary}[Not through accuracy]\label{c:acc}
The art of a photograph does not consist in its accuracy.
\end{corollary}

\begin{proof}
Accuracy, the correspondence between what a photograph shows and how the scene looked from the camera's position, is guaranteed by transparency for every photograph made by an unmanipulated process, and so is shared by the photograph that is a work of art and by the image of a fixed surveillance camera trained on the same scene. A property shared by all photographs of a scene cannot be what distinguishes those that are artworks, nor fix the aesthetic properties that, by Lemma~\ref{l:aes}, distinguish one photograph of the scene from another.
\end{proof}

Corollary~\ref{c:acc} is the point on which AES-16 relies: representational accuracy is not the measure of pictorial value, since if it were, every photograph would surpass every painting. Theorem~\ref{t:main} also shows what is right in Scruton's argument. A photograph does not represent its subject through the photographer's thought, as a painting does; what the photographer's thought controls is which seeing of the subject the photograph offers. The photograph is a presented seeing: a view of the world, transparent to it, that someone chose and offered as worth looking at.

\section{Objections and replies}

\begin{objection}[Interest in the subject]\label{o:subj}
To take an interest in a photograph is to take an interest in what it shows. If the photograph is transparent, the interest is in the subject, seen through the photograph, and the photograph itself is not an object of aesthetic interest \citep{Scruton1981}.
\end{objection}

\begin{reply}
The interest is in the subject as seen from this viewpoint, at this moment, within this frame. That is an interest in a way of seeing the subject, and the way of seeing is not a property of the subject. The objection runs together two things transparency keeps apart: what is seen, which the scene fixes, and how it is seen, which framing fixes. Looking at Cartier-Bresson's photograph of a man leaping across a puddle behind the Gare Saint-Lazare is not looking at the man as one might have seen him on the day; no one on the day saw him in that frame, with his reflection completing the figure, at that instant. The interest in the photograph is an interest in what it shows, but what it shows includes its composition.
\end{reply}

\begin{objection}[Thin agency]\label{o:thin}
Framing and selection are thin. The photographer does not make the forms in the frame; she finds them. Finding is not making, and what is merely found cannot be the photographer's achievement.
\end{objection}

\begin{reply}
Finding, in the relevant sense, is choosing which view to make, and it is an exercise of competence that can succeed or fail, which is what Lemma~\ref{l:attr} requires. Selection is a recognized form of artistic agency elsewhere, in the readymade, in collage and in the editing of film, and on the account of AES-01 a work's status as art does not depend on the maker's having shaped its matter by hand. The companion manuscript AES-14-b develops the parallel. The objection would prove too much: the painter of a landscape also does not make the forms of the hills, and her achievement lies partly in what she chose to show and from where.
\end{reply}

\begin{objection}[Automation]\label{o:auto}
Contemporary cameras choose focus, exposure and tonal rendering by algorithm, and some choose the moment. If the settings are not chosen by the photographer, framing is not hers either.
\end{objection}

\begin{reply}
Automation transfers some of the choices of Definition~\ref{d:choices} to the device, and to that extent reduces the scope of the photographer's achievement. It does not remove viewpoint, frame or selection, and these suffice for Lemmas~\ref{l:aes} and~\ref{l:attr}. A photograph whose framing and selection were also automated, such as the output of a fixed camera chosen by no one, has aesthetic properties but no maker's achievement in them, and Theorem~\ref{t:main} does not make it a work of art. This is the right result: the art of photography lies where the choices lie.
\end{reply}

\begin{objection}[Manipulation]\label{o:manip}
The photographs most often admired as art are staged, retouched or composited. Their art lies in their departures from transparency, as Scruton said, and photography as such is not an art.
\end{objection}

\begin{reply}
Some photographic art does depend on manipulation, and such works are hybrid, as Scruton held. But the argument of Section~3 applies to unmanipulated photographs: the street photography of Cartier-Bresson, much documentary and landscape photography, and the portrait made by choosing the moment rather than by retouching it. These are transparent in Walton's sense, and they have the aesthetic properties and the achievement value that Lemmas~\ref{l:aes} and~\ref{l:attr} describe. The existence of a hybrid art does not show that the pure one is not an art.
\end{reply}

\section{What the result does not establish}

The result does not show that every photograph is a work of art; most are not, since they are made without the choices, or without the competence, that Lemma~\ref{l:attr} requires, or are not presented and taken up as art. It does not settle whether photographs represent in Scruton's sense, and is compatible with his claim that they do not. It does not deny transparency; it requires only that transparency be understood as a dependence conditional on framing. Finally, it does not rank photography among the arts or compare its achievements with those of painting.

\section{Conclusion}

Photography is art. Its causal transparency, which Walton described and Scruton took to exclude the photographer's thought from the image, is a dependence of what a photograph shows on the scene given its viewpoint, frame, moment and settings. Those are chosen, and the exposure presented is selected from others. Framing and selection fix aesthetic properties that no scene has, and they are appreciable as the photographer's achievement. The art of photography lies there, and not in the accuracy that every photograph of a scene shares.

\begin{thebibliography}{99}\small
\bibitem[Bazin(1960)]{Bazin1960} Bazin, A. (1960). The ontology of the photographic image. Trans. H. Gray. \emph{Film Quarterly}, 13(4), 4--9.
\bibitem[Cartier-Bresson(1952)]{CartierBresson1952} Cartier-Bresson, H. (1952). \emph{The Decisive Moment}. Simon and Schuster.
\bibitem[Cohen and Meskin(2004)]{CohenMeskin2004} Cohen, J. and Meskin, A. (2004). On the epistemic value of photographs. \emph{The Journal of Aesthetics and Art Criticism}, 62(2), 197--210.
\bibitem[Costello(2018)]{Costello2018} Costello, D. (2018). \emph{On Photography: A Philosophical Inquiry}. Routledge.
\bibitem[Currie(1991)]{Currie1991} Currie, G. (1991). Photography, painting and perception. \emph{The Journal of Aesthetics and Art Criticism}, 49(1), 23--29.
\bibitem[Friday(2002)]{Friday2002} Friday, J. (2002). \emph{Aesthetics and Photography}. Ashgate.
\bibitem[King(1992)]{King1992} King, W.~L. (1992). Scruton and reasons for looking at photographs. \emph{The British Journal of Aesthetics}, 32(3), 258--265.
\bibitem[Lopes(2016)]{Lopes2016} Lopes, D.~M. (2016). \emph{Four Arts of Photography: An Essay in Philosophy}. Wiley-Blackwell.
\bibitem[Scruton(1981)]{Scruton1981} Scruton, R. (1981). Photography and representation. \emph{Critical Inquiry}, 7(3), 577--603.
\bibitem[Snyder and Allen(1975)]{SnyderAllen1975} Snyder, J. and Allen, N.~W. (1975). Photography, vision, and representation. \emph{Critical Inquiry}, 2(1), 143--169.
\bibitem[Talbot(1844)]{Talbot1844} Talbot, W.~H.~F. (1844). \emph{The Pencil of Nature}. Longman, Brown, Green and Longmans.
\bibitem[Walton(1984)]{Walton1984} Walton, K.~L. (1984). Transparent pictures: On the nature of photographic realism. \emph{Critical Inquiry}, 11(2), 246--277.
\end{thebibliography}

\end{document}
